\begin{table}[t]
\centering
\caption{$\lambda$ sweep comparing \AmbGCE{} and \TanGCE{} at fixed step sizes across three concept domains (Qwen2.5-0.5B, layer~8). NL: retrained nonlinear MLP probe accuracy (hidden 256, 1500 steps, AdamW lr=$10^{-3}$; same architecture used for HPO and evaluation) on the target concept (lower = stronger erasure). Both methods use the same GCE concept scorer (hidden 256, 1500 steps). \AmbGCE{} becomes less effective at high $\lambda$ under the current protocol, while \TanGCE{} continues improving. We interpret this as supportive evidence for the local manifold motivation of tangent-constrained updates, not as a complete characterization of the leakage-utility frontier. The expanded HPO grid $\{0.25,\ldots,32.0\}$ used in the multi-model experiments further tests this range. Bold: best NL per domain.}
\label{tab:lambda_sweep}
\small
\setlength{\tabcolsep}{5pt}
\begin{tabular}{ll|ccc}
\toprule
 & & \multicolumn{3}{c}{NL$\downarrow$ at $\lambda=$} \\
Concept & Method & $0.5$ & $4.0$ & $8.0$ \\
\midrule
\multirow{2}{*}{Sycophancy}
 & \AmbGCE{} & .678 & .694 & .734 \\
 & \TanGCE{}  & .662 & .596 & \textbf{.588} \\
\midrule
\multirow{2}{*}{Gender}
 & \AmbGCE{} & .666 & .700 & .702 \\
 & \TanGCE{}  & .704 & .580 & \textbf{.470} \\
\midrule
\multirow{2}{*}{Safety}
 & \AmbGCE{} & .664 & .776 & .770 \\
 & \TanGCE{}  & .748 & .680 & \textbf{.646} \\
\bottomrule
\end{tabular}
\end{table}
